% Results table for Sec. \ref{sec:results}; \input from sections/refinements.tex
% (not experiments.tex) so that the two-column float lands at the top of the page
% where Experiments begins.
\begin{table*}[t]
\centering
\caption{Test-set Recall@5/15/25/35\% and Recall-AUC of the baseline and of each
proposed refinement (mean $\pm$ std over $100$ simulation runs). Best value per
column in bold.}
\label{tab:results}
\footnotesize
\setlength{\tabcolsep}{3pt}
\begin{tabular}{@{}lccccc@{}}
\toprule
Variant & Recall@5\% & Recall@15\% & Recall@25\% & Recall@35\% & Recall-AUC \\
\midrule
Baseline
  & $0.3588 \pm 0.0782$ & $0.8185 \pm 0.0242$ & $0.9236 \pm 0.0125$ & $0.9593 \pm 0.0116$ & $0.8914 \pm 0.0118$ \\
\midrule
Forward-pass $f_j(t)$ (\S\ref{sec:forward-f})
  & $0.4512 \pm 0.0472$ & $\mathbf{0.8519} \pm 0.0294$ & $\mathbf{0.9380} \pm 0.0190$ & $\mathbf{0.9651} \pm 0.0108$ & $0.9088 \pm 0.0114$ \\
Negative feedback, second SAE (\S\ref{sec:negative-feedback})
  & $0.4663 \pm 0.0432$ & $0.8507 \pm 0.0307$ & $0.9352 \pm 0.0202$ & $0.9647 \pm 0.0109$ & $\mathbf{0.9095} \pm 0.0110$ \\
Recency-weighted $\phi_j$ (\S\ref{sec:recency-phi})
  & $\mathbf{0.4927} \pm 0.0388$ & $0.8396 \pm 0.0251$ & $0.9306 \pm 0.0139$ & $0.9623 \pm 0.0106$ & $0.9083 \pm 0.0104$ \\
Variance-based $p_j$ (\S\ref{sec:variance-p})
  & $0.4566 \pm 0.0402$ & $0.8341 \pm 0.0263$ & $0.9292 \pm 0.0138$ & $0.9621 \pm 0.0106$ & $0.9046 \pm 0.0104$ \\
Factored score, scheduled coeff. (\S\ref{sec:factored-score})
  & $0.4430 \pm 0.0402$ & $0.8330 \pm 0.0269$ & $0.9294 \pm 0.0137$ & $0.9620 \pm 0.0106$ & $0.9033 \pm 0.0107$ \\
\bottomrule
\end{tabular}
\end{table*}


\section{Experiments}
\label{sec:experiments}

\subsection{Methodology}
\label{sec:methodology}

For each variant we grid-search its hyperparameters and select the
configuration maximizing validation-set Recall-AUC,
$\mathrm{AUC} = \int_0^1 \mathrm{Recall}(\tau)\,d\tau$, where
$\mathrm{Recall}(\tau)$ is recall at sending fraction $\tau$. Variants are
compared on test-set Recall@5/15/25/35\% and Recall-AUC, as mean $\pm$ std
over repeated simulations. Before testing modifications we reproduced the
original model on our pipeline; at the operating points reported by
\citet{zid2025active} (25/50/75\%) our reproduction matches the published
point estimates exactly ($0.923$, $0.975$, $0.989$), with larger variance
attributable to fewer simulation runs.

\subsection{Results}
\label{sec:results}

Table~\ref{tab:results} reports the test-set results of the baseline and of
each refinement, as mean $\pm$ std over $100$ simulation runs. Every refinement
improves Recall-AUC over the baseline (by $1.19$ to $1.81$ percentage points),
and the largest gains appear at the smallest sending fraction: at 5\% the
improvement in recall ranges from $8.42$ to $13.39$ percentage points, while by
35\% it is below $0.6$ percentage points for all variants.
